\documentclass[11pt,a4paper]{article}

\usepackage[T1]{fontenc}
\usepackage{lmodern}
\usepackage[margin=22mm,headheight=15pt]{geometry}
\usepackage{amsmath,amssymb,mathtools}
\usepackage{microtype}
\usepackage{xcolor}
\usepackage{enumitem}
\usepackage{fancyhdr}
\usepackage{titlesec}
\usepackage{needspace}

\definecolor{examblue}{HTML}{17365D}
\definecolor{rulegray}{HTML}{D6DCE4}

\pagestyle{fancy}
\fancyhf{}
\lhead{\textsc{Machine Learning 1}}
\rhead{\textsc{Mock Exam A}}
\cfoot{\thepage}
\renewcommand{\headrulewidth}{0.4pt}

\titleformat{\section}
  {\Large\bfseries\color{examblue}}
  {}{0pt}{}
  [\vspace{0.15em}\color{rulegray}\titlerule]
\titleformat{\subsection}
  {\normalsize\bfseries\color{examblue}}
  {}{0pt}{}

\setlength{\parindent}{0pt}
\setlength{\parskip}{0.45em}
\setlist[itemize]{leftmargin=1.5em,itemsep=0.1em,topsep=0.25em}
\setlist[enumerate]{leftmargin=1.7em,itemsep=0.65em,topsep=0.35em}
\allowdisplaybreaks

\newcommand{\points}[1]{\unskip\nobreak\hfill\mbox{\textbf{[#1 points]}}}
\newcommand{\exercise}[3]{%
  \Needspace{8\baselineskip}%
  \par\addvspace{1.1em}%
  \noindent
  \begin{minipage}[t]{0.73\textwidth}
    \vspace{0pt}\raggedright\large\bfseries\color{examblue}
    Exercise #1: #2
  \end{minipage}\hfill
  \begin{minipage}[t]{0.24\textwidth}
    \vspace{0pt}\raggedleft\small\color{examblue}#3
  \end{minipage}%
  \par\nobreak\vspace{0.2em}%
  {\color{rulegray}\hrule height 0.5pt}%
  \vspace{0.55em}%
}

\begin{document}

\begin{center}
  {\Huge\bfseries\color{examblue} Machine Learning 1}\par
  \vspace{0.35em}
  {\LARGE\bfseries Mock Exam A}\par
  \vspace{0.9em}
  {\large Duration: 2 hours \qquad Total: 100 points}
\end{center}

\vspace{0.8em}
\begin{tabular*}{\textwidth}{@{\extracolsep{\fill}}ll}
  \textbf{Name:} \rule{7cm}{0.4pt} & \textbf{Student ID:} \rule{4cm}{0.4pt}
\end{tabular*}

\vspace{1em}
\fcolorbox{rulegray}{white}{%
  \begin{minipage}{0.94\textwidth}
    \textbf{Instructions}
    \begin{itemize}
      \item Show all important intermediate steps. Answers without justification can receive reduced credit.
      \item State any additional assumptions that you use.
      \item Give matrix and vector dimensions where the question asks for them.
      \item This course uses numerator layout: if $f:\mathbb{R}^m\to\mathbb{R}^n$, then its Jacobian has shape $n\times m$.
      \item Unless a question states otherwise, vectors are column vectors.
      \item The suggested working times total 110 minutes. Reserve about 10 minutes to check your answers.
    \end{itemize}
  \end{minipage}%
}

\exercise{1}{Multivariate derivatives and regularization}{20 points, 18 minutes}

\subsection*{1.1 Temperature-scaled log-softmax}

Let $x\in\mathbb{R}^d$ and let $\tau>0$. Define $s:\mathbb{R}^d\to\mathbb{R}^d$ and $\ell:\mathbb{R}^d\to\mathbb{R}^d$ by
\begin{align*}
  s_i(x) &= \frac{\exp(x_i/\tau)}{\sum_{k=1}^{d}\exp(x_k/\tau)},
  &\ell_i(x) &= \log s_i(x).
\end{align*}

\begin{enumerate}
  \item Derive $\partial \ell_i/\partial x_j$ using index notation. Express the answer using $s_j$, $\tau$, and the Kronecker delta $\delta_{ij}$. \points{5}
  \item Write the full Jacobian $J_{\ell}(x)$ in matrix notation and state its dimensions. \points{4}
  \item Show that $J_{\ell}(x)\mathbf{1}=0$, where $\mathbf{1}$ is the all-ones vector. Explain this result using a transformation of $x$ that leaves $\ell(x)$ unchanged. \points{4}
\end{enumerate}

\subsection*{1.2 Generalized ridge regression}

Let $X\in\mathbb{R}^{n\times p}$, $y\in\mathbb{R}^n$, $\theta\in\mathbb{R}^p$, and $\gamma>0$. Let $R\in\mathbb{R}^{p\times p}$ be symmetric positive definite. Consider
\[
  L(\theta)=\frac{1}{n}\lVert X\theta-y\rVert_2^2+\gamma\theta^\top R\theta.
\]

\begin{enumerate}[resume]
  \item Derive the stationary point of $L(\theta)$. Then show that it is the unique global minimizer, even when $X$ does not have full column rank. \points{7}
\end{enumerate}

\exercise{2}{Waiting times and Bayesian estimation}{24 points, 25 minutes}

A monitoring system checks a component repeatedly. Each check succeeds independently with probability $q$, where $0<q<1$. Let $K$ be the number of failed checks before the first success. Its probability mass function is
\[
  p(K=k\mid q)=q(1-q)^k,\qquad k=0,1,2,\ldots
\]
You may use
\[
  \sum_{k=0}^{\infty}r^k=\frac{1}{1-r},\qquad |r|<1.
\]

\begin{enumerate}
  \item Derive $\mathbb{E}[K\mid q]$. Do not only state the result. \points{3}
  \item Compute $P(K=2\mid q=0.25)$. Give an exact or decimal value. \points{3}
\end{enumerate}

Now suppose that $K_1,\ldots,K_N$ are independent observations and define $S=\sum_{n=1}^{N}K_n$. Assume $S>0$.

\begin{enumerate}[resume]
  \item Write the likelihood and log-likelihood of $q$, up to constants that do not depend on $q$. \points{3}
  \item Derive the maximum-likelihood estimator of $q$. Verify that your stationary point is a maximum. \points{4}
\end{enumerate}

Place a $\operatorname{Beta}(a,b)$ prior on $q$, with density
\[
  p(q)\propto q^{a-1}(1-q)^{b-1},\qquad 0<q<1,
\]
where $a>1$ and $b>1$.

\begin{enumerate}[resume]
  \item Derive the posterior distribution of $q$ and identify its parameters. \points{4}
  \item Derive the MAP estimator of $q$. \points{4}
  \item Identify a one-dimensional statistic of the data that contains all data dependence in the posterior. Briefly justify your answer. \points{3}
\end{enumerate}

\exercise{3}{Estimating an unknown range}{14 points, 14 minutes}

Measurements $X_1,\ldots,X_N$ are independent and uniformly distributed on $[-M,M]$, where $M>0$ is unknown. Define
\[
  A=\max_{1\leq n\leq N}|X_n|.
\]
Assume that at least one observed value is nonzero, so $A>0$.

\begin{enumerate}
  \item Write the likelihood $p(x_1,\ldots,x_N\mid M)$, including the support constraint. \points{4}
  \item Derive the maximum-likelihood estimator of $M$. For the observations
  \[
    -3.2,\quad 1.1,\quad 2.4,\quad -2.8,\quad 0.5,
  \]
  give its numerical value. \points{3}
\end{enumerate}

You may use
\[
  \mathbb{E}[A]=\frac{NM}{N+1},
  \qquad
  \mathbb{E}[A^2]=\frac{NM^2}{N+2}.
\]

\begin{enumerate}[resume]
  \item Derive a multiplicative correction of the maximum-likelihood estimator that is unbiased for $M$. \points{3}
  \item Derive the variance of your unbiased estimator and simplify the result. \points{4}
\end{enumerate}

\exercise{4}{Multi-output regression with correlated noise}{23 points, 29 minutes}

For each observation, let $\phi_n\in\mathbb{R}^m$ be a feature vector and $t_n\in\mathbb{R}^k$ be a target vector. Let $W\in\mathbb{R}^{m\times k}$. The model is
\[
  t_n=W^\top\phi_n+\varepsilon_n,
  \qquad
  \varepsilon_n\sim\mathcal{N}(0,\sigma^2C),
\]
where $C\in\mathbb{R}^{k\times k}$ is known and symmetric positive definite, while $\sigma^2>0$ is unknown. The noise vectors are independent.

Let $\Phi\in\mathbb{R}^{N\times m}$ have rows $\phi_n^\top$, and let $T\in\mathbb{R}^{N\times k}$ have rows $t_n^\top$. Assume that $\Phi$ has full column rank, $N>m$, and the data cannot be fitted exactly.

\begin{enumerate}
  \item State the dimensions of $\Phi W$, $T-\Phi W$, and $(T-\Phi W)^\top(T-\Phi W)$. \points{2}
  \item Derive the log-likelihood of $W$ and $\sigma^2$, including all terms that depend on either parameter. You may express the quadratic term using a trace. \points{6}
  \item Derive the maximum-likelihood estimator of $W$. Show why it does not depend on $C$ or $\sigma^2$. \points{7}
  \item Using the fitted value of $W$, derive the maximum-likelihood estimator of $\sigma^2$. \points{5}
  \item Let $E=T-\Phi W_{\mathrm{ML}}$. Show that $\Phi^\top E=0$ and explain the geometric meaning of this identity. \points{3}
\end{enumerate}

\exercise{5}{Information from one Bayesian observation}{19 points, 24 minutes}

After observing a data set $\mathcal{D}_N$, the posterior over the regression weights is
\[
  p(w\mid\mathcal{D}_N)=\mathcal{N}(w\mid\mu_N,\Sigma_N),
\]
where $w,\mu_N\in\mathbb{R}^m$ and $\Sigma_N\in\mathbb{R}^{m\times m}$ is symmetric positive definite. A new observation has feature vector $\phi\in\mathbb{R}^m$, target $t\in\mathbb{R}$, and likelihood
\[
  p(t\mid w,\phi)=\mathcal{N}(t\mid w^\top\phi,\beta^{-1}),
\]
where $\beta>0$ is the observation-noise precision.

The standard precision and natural-parameter updates are
\begin{align*}
  \Sigma_{N+1}^{-1} &= \Sigma_N^{-1}+\beta\phi\phi^\top,\\
  \Sigma_{N+1}^{-1}\mu_{N+1} &= \Sigma_N^{-1}\mu_N+\beta\phi t.
\end{align*}
You may use the Sherman--Morrison identity
\[
  (A+uv^\top)^{-1}
  =A^{-1}-\frac{A^{-1}uv^\top A^{-1}}{1+v^\top A^{-1}u}.
\]

\begin{enumerate}
  \item Derive the rank-one covariance update
  \[
    \Sigma_{N+1}=\Sigma_N-
    \frac{\Sigma_N\phi\phi^\top\Sigma_N}
    {\beta^{-1}+\phi^\top\Sigma_N\phi}.
  \]
  \points{4}
  \item Derive the innovation form of the mean update
  \[
    \mu_{N+1}=\mu_N+K(t-\phi^\top\mu_N),
  \]
  and give the gain vector $K$. Explain the roles of the direction and the scalar residual in this update. \points{5}
  \item Before observing $t$, derive the posterior predictive distribution $p(t\mid\phi,\mathcal{D}_N)$. Give both its mean and variance, and identify the two sources of predictive variance. \points{4}
  \item For any direction $a\in\mathbb{R}^m$, derive
  \[
    a^\top(\Sigma_N-\Sigma_{N+1})a.
  \]
  Show that it is nonnegative. State exactly when it is zero and explain what this says about the parameter directions learned from the observation. \points{6}
\end{enumerate}

\vfill
\begin{center}
  \textbf{End of Mock Exam A}
\end{center}

\end{document}
